\frfilename{mt233.tex}
\versiondate{16.6.02}
\copyrightdate{1994}

\def\chaptername{Teorema Radon-Nikod\'ym}
\def\sectionname{Ekspektasi bersyarat}

\newsection{233}

Saya menyediakan satu bagian untuk tinjauan pertama atas salah satu
penerapan utama Teorema Radon-Nikod\'ym. Gagasan ini merupakan salah satu
gagasan terpenting dalam teori ukuran, dan akan muncul berulang kali dalam
berbagai bentuk. Di sini saya memberikan definisi dan sifat-sifat paling
dasar ekspektasi bersyarat sebagaimana muncul dalam teori probabilitas
abstrak, disertai catatan tentang fungsi konveks dan suatu versi
Ketaksamaan Jensen (233I-233J).

\leader{233A}{Subaljabar-$\sigma$} Misalkan $X$ suatu himpunan dan $\Sigma$
suatu aljabar-$\sigma$ dari subhimpunan-subhimpunan $X$. Suatu
{\bf subaljabar-$\sigma$} dari $\Sigma$ adalah aljabar-$\sigma$ $\Tau$
dari subhimpunan-subhimpunan $X$ sedemikian sehingga
$\Tau\subseteq\Sigma$. Jika $(X,\Sigma,\mu)$ suatu ruang ukur dan
$\Tau$ suatu subaljabar-$\sigma$ dari $\Sigma$, maka
$(X,\Tau,\mu\restrp\Tau)$ juga merupakan ruang ukur\cmmnt{; hal ini
langsung mengikuti definisi (112A). Sekarang kita mempunyai lema
sederhana berikut. Lema ini merupakan kasus khusus 235G di bawah,
tetapi saya memberikan bukti tersendiri apabila Anda belum hendak
memulai penyelidikan umum yang ditempuh dalam \S235}.

\leader{233B}{Lema} Misalkan $(X,\Sigma,\mu)$ suatu ruang ukur dan $\Tau$
suatu subaljabar-$\sigma$ dari $\Sigma$. Suatu fungsi bernilai real $f$
yang didefinisikan pada suatu subhimpunan dari $X$
bersifat $\mu\restrp\Tau$-terintegralkan jika dan hanya jika (i) fungsi
tersebut $\mu$-terintegralkan (ii) $\dom f$ koterabaikan terhadap
$\mu\restrp\Tau$ (iii) $f$ terukur secara virtual terhadap
$\mu\restrp\Tau$; dan dalam hal ini
$\int fd(\mu\restrp\Tau)=\int fd\mu$.

\proof{{\bf (a)} Pertama-tama perhatikan bahwa jika $f$ suatu fungsi
sederhana terhadap $\mu\restrp\Tau$, yakni dapat dinyatakan sebagai
$\sum_{i=0}^na_i\chi E_i$ dengan $a_i\in\Bbb R$, $E_i\in\Tau$, dan
$(\mu\restrp\Tau)E_i<\infty$ untuk setiap $i$, maka $f$ merupakan fungsi
sederhana terhadap $\mu$ dan

\Centerline{$\int f d\mu=\sum_{i=0}^na_i\mu E_i
=\int f d(\mu\restrp\Tau)$.}

\medskip

{\bf (b)} Misalkan $U_{\mu}$ himpunan fungsi nonnegatif yang
$\mu$-terintegralkan dan $U_{\mu\restrp\Tau}$ himpunan fungsi nonnegatif
yang $\mu\restrp\Tau$-terintegralkan.

Andaikan $f\in U_{\mu\restrp\Tau}$. Maka terdapat barisan tak menurun
$\sequencen{f_n}$ dari fungsi-fungsi sederhana terhadap
$\mu\restrp\Tau$ sedemikian sehingga
$f(x)=\lim_{n\to\infty}f_n\,\mu\restrp\Tau$-hampir di mana-mana\ dan

\Centerline{$\int fd(\mu\restrp\Tau)
=\lim_{n\to\infty}\int f_nd(\mu\restrp\Tau)$.}

\noindent Namun kini setiap $f_n$ juga sederhana terhadap $\mu$, dan
$\int f_nd\mu=\int f_nd(\mu\restrp\Tau)$ untuk setiap $n$, serta
$f=\lim_{n\to\infty}f_n\,\mu$-hampir di mana-mana. Jadi
$f\in U_{\mu}$ dan $\int fd\mu=\int fd(\mu\restrp\Tau)$.

\medskip

{\bf (c)} Sekarang andaikan $f$ bersifat
$\mu\restrp\Tau$-terintegralkan. Maka fungsi tersebut merupakan selisih dua anggota
$U_{\mu\restrp\Tau}$, sehingga $\mu$-terintegralkan, dan
$\int fd\mu=\int fd(\mu\restrp\Tau)$. Syarat (ii) dan (iii) juga
terpenuhi, sesuai dengan konvensi yang ditetapkan dalam Jilid 1
(122Nc, 122P-122Q).

\medskip

{\bf (d)} Andaikan $f$ memenuhi syarat (i)-(iii). Maka
$|f|\in U_{\mu}$, dan terdapat himpunan koterabaikan
$E\subseteq\dom f$ sedemikian sehingga $E\in\Tau$ dan $f\restr E$
terukur terhadap $\Tau$. Oleh karena itu, $|f|\restr E$ terukur terhadap
$\Tau$. Sekarang, jika $\epsilon>0$, maka

\Centerline{$(\mu\restrp\Tau)\{x:x\in E,\,|f|(x)\ge\epsilon\}
=\mu\{x:x\in E,\,|f|(x)\ge\epsilon\}
\le{1\over\epsilon}\int|f|d\mu<\infty$;}

\noindent selain itu,

$$\eqalign{\sup\{\int g\,d(\mu\restrp\Tau):g&\text{ adalah suatu fungsi }
  \mu\restrp\Tau\text{-sederhana},
  \,g\le |f|\,\mu\restrp\Tau\text{-hampir di mana-mana}\}\cr
&=\sup\{\int g\,d\mu:g\text{ adalah suatu fungsi }
  \mu\restrp\Tau\text{-sederhana},
  \,g\le |f|\,\mu\restrp\Tau\text{-hampir di mana-mana}\}\cr
&\le\sup\{\int g\,d\mu:g\text{ adalah suatu fungsi }
  \mu\text{-sederhana},\,g\le |f|\,\mu\text{-hampir di mana-mana}\}\cr
&\le\int |f|d\mu<\infty.\cr}$$

\noindent Menurut kriteria 122Ja, $|f|\in U_{\mu\restrp\Tau}$.
Akibatnya, karena $f$ terukur secara virtual terhadap $\mu\restrp\Tau$
sebagai fungsi yang terukur terhadap $\Tau$, fungsi tersebut
$\mu\restrp\Tau$-terintegralkan menurut 122P. Ini menuntaskan bukti.
}%end of proof of 233B

\cmmnt{
\leader{233C}{Catatan (a)} Argumen tepat di atas saya uraikan dengan
sangat terperinci, bahkan hingga terkesan bertele-tele. Namun, walaupun
saya dapat dituduh membuang-buang kertas dengan menuliskan semuanya,
menurut saya setiap unsur argumen itu diperlukan bagi hasilnya. Memang,
sebagian perincian diperlukan hanya karena saya menggunakan pengertian
`fungsi terintegralkan` yang begitu luas; jika Anda membatasi pengertian
`keterintegralan` pada fungsi terukur yang didefinisikan pada seluruh
ruang ukur, pada tahap ini terdapat beberapa penyederhanaan, yang kelak
harus dibayar ketika Anda mendapati bahwa banyak penerapan utama
menyangkut fungsi yang didefinisikan oleh rumus-rumus yang tidak berlaku
pada seluruh ruang yang mendasarinya.

Pokok penting yang harus dipahami ialah bahwa meskipun himpunan yang
terabaikan terhadap $\mu\restrp\Tau$ selalu terabaikan terhadap $\mu$,
himpunan yang terabaikan terhadap $\mu$ belum tentu terabaikan terhadap
$\mu\restrp\Tau$.

\header{233Cb}{\bf (b)} Sebagai contoh sesederhana mungkin mengenai
masalah yang dapat timbul, saya menawarkan contoh berikut. Misalkan
$(X,\Sigma,\mu)$ adalah $[0,1]^2$ dengan ukuran Lebesgue. Misalkan
$\Tau$ himpunan anggota $\Sigma$ yang dapat dinyatakan sebagai
$F\times[0,1]$ untuk suatu $F\subseteq[0,1]$; mudah dilihat bahwa $\Tau$
merupakan subaljabar-$\sigma$ dari $\Sigma$. Tinjau $f$,
$g:X\to[0,1]$ yang didefinisikan dengan

\Centerline{$f(t,u)=1$ jika $u>0$, $0$ selainnya,}

\Centerline{$g(t,u)=1$ jika $t>0$, $0$ selainnya.}

\noindent Maka $f$ dan $g$ keduanya $\mu$-terintegralkan karena konstan
$\mu$-hampir di mana-mana. Namun, hanya $g$ yang
$\mu\restrp\Tau$-terintegralkan, sebab setiap $E\in\Tau$ yang tidak
terabaikan memuat suatu irisan vertikal lengkap
$\{t\}\times[0,1]$, sehingga $f$ mengambil kedua nilai $0$ dan $1$ pada
$E$. Jika kita menetapkan

\Centerline{$h(t,u)=1$ jika $u>0$, tidak terdefinisi selainnya,}

\noindent maka sekali lagi (menurut konvensi yang saya gunakan) $h$
$\mu$-terintegralkan tetapi tidak $\mu\restrp\Tau$-terintegralkan, karena
tidak ada anggota koterabaikan dari $\Tau$ yang termuat dalam domain $h$.

\header{233Cc}{\bf (c)} Jika $f$ didefinisikan pada seluruh $X$ dan
$\mu\restrp\Tau$ lengkap, maka tentu saja $f$
$\mu\restrp\Tau$-terintegralkan jika dan hanya jika fungsi tersebut
$\mu$-terintegralkan dan terukur terhadap $\Tau$. Namun, perhatikan bahwa
dalam contoh tepat di atas, yang merupakan salah satu pola dasar topik
ini, $\mu\restrp\Tau$ tidak lengkap, karena himpunan-himpunan satu titik
terabaikan tetapi tidak terukur.
}%end of comment

\leader{233D}{Ekspektasi bersyarat} Misalkan $(X,\Sigma,\mu)$ suatu
ruang probabilitas\cmmnt{, yakni ruang ukur dengan $\mu X=1$}.
\cmmnt{(Hampir semua gagasan di sini berlaku dengan sempurna bagi
sembarang ruang ukur berhingga total, tetapi tampaknya tidak ada manfaat
yang diperoleh dari perluasan tersebut, dan ungkapan tradisional
`ekspektasi bersyarat` menuntut ruang probabilitas.)}
Misalkan $\Tau\subseteq\Sigma$ suatu subaljabar-$\sigma$.
\dvro{ Jika $f$ suatu fungsi bernilai real yang $\mu$-terintegralkan,
terdapat fungsi $\mu\restrp\Tau$-terintegralkan $g$ sedemikian
sehingga $\int_Fg=\int_Ff$ untuk setiap $F\in\Tau$; fungsi semacam itu
merupakan {\bf suatu ekspektasi bersyarat} dari $f$ pada $\Tau$.}{}

\cmmnt{\header{233Da}{\bf (a)} Untuk setiap fungsi bernilai real yang
$\mu$-terintegralkan, katakanlah $f$, dan didefinisikan pada suatu subhimpunan
koterabaikan dari $X$, kita mempunyai integral tak tentu yang
bersesuaian, $\nu_f:\Sigma\to\Bbb R$, yang diberikan oleh rumus
$\nu_fE=\int_Ef$ untuk setiap $E\in\Sigma$. Kita tahu bahwa $\nu_f$
aditif terhitung dan kontinu sejati terhadap $\mu$, yang dalam konteks
sekarang sama dengan mengatakan bahwa fungsional itu kontinu mutlak
(232Bc-232Bd). Sekarang tinjau pembatasan $\mu\restrp\Tau$ dan
$\nu_f\restrp\Tau$ dari $\mu$ dan $\nu_f$ pada aljabar-$\sigma$ $\Tau$.
Langsung dari definisi `aditif terhitung` dan `kontinu mutlak` diperoleh
bahwa $\nu_f\restrp\Tau$ aditif terhitung dan kontinu mutlak terhadap
$\mu\restrp\Tau$, sehingga kontinu sejati terhadap
$\mu\restrp\Tau$. Akibatnya, Teorema Radon-Nikod\'ym (232E) menyatakan
bahwa terdapat fungsi $\mu\restrp\Tau$-terintegralkan $g$ sedemikian
sehingga $(\nu_f\restrp\Tau)F=\int_Fg\,d(\mu\restrp\Tau)$ untuk setiap
$F\in\Tau$.
}%end of comment

\cmmnt{\header{233Db}{\bf (b)} Mari kita definisikan {\bf suatu
ekspektasi bersyarat dari $f$ pada} $\Tau$ sebagai fungsi semacam itu;
yakni, suatu fungsi $\mu\restrp\Tau$-terintegralkan $g$ sedemikian
sehingga $\int_Fg\,d(\mu\restrp\Tau)=\int_Ffd\mu$ untuk setiap
$F\in\Tau$. Dengan menengok kembali 233B, kita melihat bahwa untuk $g$
semacam itu berlaku

\Centerline{$\int_Fg\,d(\mu\restrp\Tau)
=\int g\times\chi F\,d(\mu\restrp\Tau)
=\int g\times\chi F\, d\mu=\int_Fg\,d\mu$}

\noindent untuk setiap $F\in\Tau$; selain itu, $g$ hampir di mana-mana
sama dengan suatu fungsi terukur terhadap $\Tau$ yang didefinisikan pada
seluruh $X$ dan juga merupakan ekspektasi bersyarat dari $f$ pada
$\Tau$ (232He).
}%end of comment

\cmmnt{\header{233Dc}{\bf (c)} Saya mencetak tebal kata `suatu` dalam
ungkapan `suatu ekspektasi bersyarat` untuk menekankan bahwa fungsi $g$
sama sekali tidak unik. Dalam 242J saya akan kembali ke pokok ini dan
menjelaskan suatu objek yang dengan tepat dapat disebut `ekspektasi
bersyarat` dari $f$ pada $\Tau$. $g$ `unik secara esensial` hanya dalam
arti bahwa jika $g_1$ dan $g_2$ keduanya merupakan ekspektasi bersyarat
dari $f$ pada $\Tau$, maka
$g_1=g_2\,\,\mu\restrp\Tau$-hampir di mana-mana\
(131Hb). Tentu saja,
hal ini berarti bahwa sangat banyak sifatnya -- misalnya fungsi
distribusi $G(a)=\hat\mu\{x:g(x)\le a\}$, dengan $\hat\mu$ pelengkapan
$\mu$ (212C) -- tidak bergantung pada $g$ mana yang kita pilih.
}%end of comment

\cmmnt{\header{233Dd}{\bf (d)} Ungkapan `ekspektasi bersyarat` perlu
dijelaskan. Ungkapan ini berasal dari identifikasi baku probabilitas
dengan ukuran, berkat Kolmogorov, yang akan saya bahas lebih lengkap
dalam Bab 27. Suatu peubah acak bernilai real dapat dipandang sebagai
fungsi terukur, atau terukur secara virtual, $f$ pada ruang probabilitas
$(X,\Sigma,\mu)$; `ekspektasinya` diidentifikasi dengan $\int f d\mu$,
dengan mengandaikan integral ini ada. Jika $F\in\Sigma$ dan $\mu F>0$,
maka `ekspektasi bersyarat dari $f$ apabila $F$ diketahui` adalah
${1\over{\mu F}}\int_Ff$. Jika $F_0,\ldots,F_n$ merupakan partisi $X$
menjadi himpunan-himpunan terukur yang berukuran tak nol, maka fungsi
$g$ yang diberikan oleh

\Centerline{$g(x)=\Bover1{\mu F_i}\int_{F_i}f$ jika $x\in F_i$}

\noindent merupakan semacam ekspektasi bersyarat yang telah
diantisipasi; jika kelak kita diberi tahu bahwa $x\in F_i$, maka $g(x)$
akan menjadi taksiran kita selanjutnya atas ekspektasi $f$. Menurut
definisi di atas, $g$ merupakan ekspektasi bersyarat dari $f$ pada
aljabar hingga $\Tau$ yang dibangkitkan oleh
$\{F_0,\ldots,F_n\}$. Intuisi yang sesuai bagi aljabar-$\sigma$ umum
$\Tau$ ialah bahwa aljabar tersebut terdiri atas kejadian-kejadian yang
akan dapat kita amati pada suatu waktu mendatang $t_0$ yang ditetapkan,
sedangkan seluruh aljabar $\Sigma$ terdiri atas semua kejadian, termasuk
yang baru dapat diamati setelah $t_0$, jika memang pernah dapat diamati.
}%end of comment

\leader{233E}{}\cmmnt{ Saya mencantumkan beberapa fakta elementer
mengenai ekspektasi bersyarat.

\medskip

\noindent}{\bf Proposisi} Misalkan $(X,\Sigma,\mu)$ suatu ruang
probabilitas dan $\Tau$ suatu subaljabar-$\sigma$ dari $\Sigma$.
Misalkan $\sequencen{f_n}$ suatu barisan fungsi bernilai real yang
$\mu$-terintegralkan, dan untuk setiap $n$ misalkan $g_n$ suatu
ekspektasi bersyarat dari $f_n$ pada $\Tau$. Maka

(a) $g_1+g_2$ merupakan ekspektasi bersyarat dari $f_1+f_2$ pada $\Tau$;

(b) untuk setiap $c\in\Bbb R$, $cg_0$ merupakan ekspektasi bersyarat dari
$cf_0$ pada $\Tau$;

(c) jika $f_1\leae f_2$, maka $g_1\leae g_2$;

(d) jika $\sequencen{f_n}$ tak menurun hampir di mana-mana\ dan
$f=\lim_{n\to\infty}f_n$ bersifat $\mu$-terintegralkan, maka
$\lim_{n\to\infty}g_n$ merupakan ekspektasi bersyarat dari $f$ pada
$\Tau$;

(e) jika $f=\lim_{n\to\infty}f_n$ terdefinisi hampir di mana-mana\ dan
terdapat fungsi $\mu$-terintegralkan $h$ sedemikian sehingga
$|f_n|\leae h$ untuk setiap $n$, maka $\lim_{n\to\infty}g_n$ merupakan
ekspektasi bersyarat dari $f$ pada $\Tau$;

(f) jika $F\in\Tau$, maka $g_0\times \chi F$ merupakan ekspektasi
bersyarat dari $f_0\times\chi F$ pada $\Tau$;

(g) jika $h$ suatu fungsi bernilai real berbatas yang terukur secara
virtual terhadap $\mu\restrp\Tau$ dan didefinisikan
$\mu\restrp\Tau$-hampir di mana-mana dalam $X$, maka $g_0\times h$
merupakan ekspektasi bersyarat dari $f_0\times h$ pada $\Tau$;

(h) jika $\Upsilon$ suatu subaljabar-$\sigma$ dari $\Tau$, maka suatu
fungsi $h_0$ merupakan ekspektasi bersyarat dari $f_0$ pada $\Upsilon$
jika dan hanya jika fungsi tersebut merupakan ekspektasi bersyarat dari
$g_0$ pada $\Upsilon$.

\proof{{\bf (a)-(b)} Kita hanya perlu mengamati bahwa

\Centerline{$\int_Fg_1+g_2d(\mu\restrp\Tau)
=\int_Fg_1d(\mu\restrp\Tau)+\int_Fg_2d(\mu\restrp\Tau)
=\int_Ff_1d\mu+\int_Ff_2d\mu=\int_Ff_1+f_2d\mu$,}

\Centerline{$\int_Fcg_0d(\mu\restrp\Tau)
=c\int_Fg_0d(\mu\restrp\Tau)=c\int_Ff_0d\mu=\int_Fcf_0d\mu$}

\noindent untuk setiap $F\in\Tau$.

\medskip

{\bf (c)} Jika $F\in\Tau$, maka

\Centerline{$\int_Fg_1d(\mu\restrp\Tau)=\int_Ff_1d\mu\le\int_Ff_2d
\mu
=\int_Fg_2d(\mu\restrp\Tau)$}

\noindent untuk setiap $F\in\Tau$; akibatnya
$g_1\le g_2\,\mu\restrp\Tau$-hampir di mana-mana\ (131Ha).

\medskip

{\bf (d)} Menurut (c), $\sequencen{g_n}$ tak menurun
$\mu\restrp\Tau$-hampir di mana-mana; selain itu,

\Centerline{$\sup_{n\in\Bbb N}\int g_nd(\mu\restrp\Tau)
=\sup_{n\in\Bbb N}\int f_nd\mu
=\int fd\mu<\infty$.}

\noindent Menurut Teorema B.Levi, $g=\lim_{n\to\infty}g_n$ terdefinisi
$\mu\restrp\Tau$-hampir di mana-mana, dan

\Centerline{$\int_Fg\,d(\mu\restrp\Tau)
=\lim_{n\to\infty}\int_Fg_nd(\mu\restrp\Tau)
=\lim_{n\to\infty}\int_Ff_nd\mu
=\int_Ffd\mu$}

\noindent untuk setiap $F\in\Tau$, sehingga $g$ merupakan ekspektasi
bersyarat dari $f$ pada $\Tau$.

\medskip

{\bf (e)} Tetapkan $f'_n=\inf_{m\ge n}f_m$ dan
$f''_n=\sup_{m\ge n}f_m$ untuk setiap $n\in\Bbb N$. Maka kita mempunyai

\Centerline{$-h\leae f'_n\le f_n\le f''_n\leae h$,}

\noindent dan $\sequencen{f'_n}$ serta $\sequencen{f''_n}$ merupakan
barisan fungsi yang monoton hampir di mana-mana dan keduanya konvergen
hampir di mana-mana ke $f$. Untuk setiap $n$, misalkan $g'_n$ dan
$g''_n$ ekspektasi bersyarat dari $f'_n$ dan $f''_n$ pada $\Tau$.
Menurut (c) dan (d), $\sequencen{g'_n}$ dan $\sequencen{g''_n}$
merupakan barisan yang monoton hampir di mana-mana dan konvergen hampir
di mana-mana ke ekspektasi bersyarat $g'$ dan $g''$ dari $f$. Tentu saja
$g'=g''\,\,\mu\restrp\Tau$-hampir di mana-mana\ (233Dc). Selain itu,
untuk setiap $n$, $g'_n\leae g_n\leae g''_n$, sehingga
$\sequencen{g_n}$ konvergen ke
$g'\,\,\mu\restrp\Tau$-hampir di mana-mana, dan
$g=\lim_{n\to\infty}g_n$ terdefinisi hampir di mana-mana serta merupakan
ekspektasi bersyarat dari $f$ pada $\Tau$.

\medskip

{\bf (f)} Untuk setiap $H\in\Tau$,

\Centerline{$\int_Hg_0\times\chi F\,d(\mu\restrp\Tau)
=\int_{H\cap F}g_0d(\mu\restrp\Tau)
=\int_{H\cap F}f_0d\mu
=\int_Hf_0\times\chi F\,d\mu$.}

\medskip

{\bf (g)}(i) Jika $h$ benar-benar sederhana terhadap
$(\mu\restrp\Tau)$, katakanlah
$h=\sum_{i=0}^na_i\chi F_i$ dengan $F_i\in\Tau$ untuk setiap $i$, maka

\Centerline{$\int_Fg_0\times h\,d(\mu\restrp\Tau)
=\sum_{i=0}^na_i\int_Fg_0\times\chi F_id(\mu\restrp\Tau)
=\sum_{i=0}^na_i\int_Ff_0\times\chi F_i\,d\mu
=\int_Ff_0\times h\,d\mu$}

\noindent untuk setiap $F\in\Tau$. (ii) Untuk kasus umum, jika $h$
terukur secara virtual terhadap $\mu\restrp\Tau$ dan
$|h(x)|\le M\,\,\mu\restrp\Tau$-hampir di mana-mana, maka terdapat
barisan $\sequencen{h_n}$ dari fungsi-fungsi sederhana terhadap
$\mu\restrp\Tau$ yang konvergen ke $h$ hampir di mana-mana, dengan
$|h_n(x)|\le M$ untuk setiap $x$ dan $n$. Sekarang
$f_0\times h_n\to f_0\times h$ hampir di mana-mana\ dan
$|f_0\times h_n|\leae M|f_0|$ untuk setiap $n$, sedangkan
$g_0\times h_n$ merupakan ekspektasi bersyarat dari $f_0\times h_n$
untuk setiap $n$. Jadi, menurut (e),
$\lim_{n\to\infty}g_0\times h_n$ merupakan ekspektasi bersyarat dari
$f_0\times h$; tetapi fungsi ini sama hampir di mana-mana dengan
$g_0\times h$.

\medskip

{\bf (h)} Kita hanya perlu memperhatikan bahwa
$\int_Hg_0d(\mu\restrp\Tau)=\int_Hf_0d\mu$ untuk setiap
$H\in\Upsilon$, sehingga

$$\eqalign{\int_Hh_0&d(\mu\restrp\Upsilon)=\int_Hg_0d(\mu\restrp\Tau)
\text{ untuk setiap }H\in\Upsilon\cr
&\iff\int_Hh_0d(\mu\restrp\Upsilon)=\int_Hf_0d\mu
\text{ untuk setiap }H\in\Upsilon.\cr}$$
}%end of proof of 233E

\cmmnt{
\leader{233F}{Catatan} Tentu saja, masing-masing hasil di atas hampir
sepele (meskipun saya kira (e) dan (g) dapat membuat Anda berhenti
sejenak untuk berpikir jika diberikan tanpa persiapan sebelumnya).
Secara keseluruhan, hasil-hasil tersebut memberi sejumlah sifat yang
cukup kuat. Dalam \S242 saya akan menyatakannya kembali dengan bahasa
yang secara sintaksis lebih langsung, tetapi bertumpu pada tingkat
abstraksi yang lebih dalam.

Sebagai ilustrasi tentang daya ekspektasi bersyarat untuk mengejutkan
kita, saya menawarkan proposisi berikut, yang bergantung pada konsep
fungsi `konveks`.
}%end of comment

\leader{233G}{Fungsi konveks} Ingat bahwa suatu fungsi bernilai real
$\phi$ yang didefinisikan pada selang $I\subseteq\Bbb R$ disebut
{\bf konveks} jika

\Centerline{$\phi(tb+(1-t)c)\le t\phi(b)+(1-t)\phi(c)$}

\noindent setiap kali $b$, $c\in I$, dan $t\in[0,1]$.

\cmmnt{\medskip

\noindent{\bf Contoh} Rumus-rumus $|x|$, $x^2$, dan
$e^{\pm x}\pm x$ mendefinisikan fungsi konveks pada $\Bbb R$; pada
$\ooint{-1,1}$ kita mempunyai $1/(1-x^2)$; pada
$\ooint{0,\infty}$ kita mempunyai $1/x$ dan $x\ln x$; pada $[0,1]$
kita mempunyai fungsi yang bernilai nol pada $\ooint{0,1}$ dan $1$ pada
$\{0,1\}$.
}%end of comment

\leader{233H}{}\cmmnt{ Teori umum fungsi konveks bersifat luas sekaligus
penting; beberapa sifatnya yang lebih menonjol saya cantumkan dalam
233Xe. Untuk saat ini, lema berikut mencakup apa yang kita perlukan.

\medskip

\noindent}{\bf Lema} Misalkan $I\subseteq\Bbb R$ suatu selang terbuka
tak kosong\cmmnt{ (berbatas atau tak berbatas)} dan
$\phi:I\to\Bbb R$ suatu fungsi konveks.

(a) Untuk setiap $a\in I$ terdapat $b\in\Bbb R$ sedemikian sehingga
$\phi(x)\ge\phi(a)+b(x-a)$ untuk setiap $x\in I$.

(b) Jika untuk setiap $q\in I\cap\Bbb Q$ kita mengambil
$b_q\in\Bbb R$ sedemikian sehingga
$\phi(x)\ge\phi(q)+b_q(x-q)$ untuk setiap $x\in I$, maka

\Centerline{$\phi(x)=\sup_{q\in I\cap\Bbb Q}\phi(q)+b_q(x-q)$}

\noindent untuk setiap $x\in I$.

(c) $\phi$ terukur Borel.

\proof{{\bf (a)} Jika $c$, $c'\in I$, dan $c<a<c'$, maka $a$ dapat
dinyatakan sebagai $dc+(1-d)c'$ untuk suatu $d\in\ooint{0,1}$, sehingga
$\phi(a)\le d\phi(c)+(1-d)\phi(c')$ dan

$$\eqalign{{\Bover{\phi(a)-\phi(c)}{a-c}}
&\le{\Bover{d\phi(c)+(1-d)\phi(c')-\phi(c)}{dc+(1-d)c'-c}}
={\Bover{(1-d)(\phi(c')-\phi(c))}{(1-d)(c'-c)}}\cr
&={\Bover{d(\phi(c')-\phi(c))}{d(c'-c)}}
={\Bover{\phi(c')-d\phi(c)-(1-d)\phi(c')}{c'-dc-(1-d)c'}}
\le{\Bover{\phi(c')-\phi(a)}{c'-a}}.\cr}$$

\noindent Ini berarti bahwa

\Centerline{$b=\sup_{c<a,c\in I}\Bover{\phi(a)-\phi(c)}{a-c}$}

\noindent berhingga, dan
$b\le\Bover{\phi(c')-\phi(a)}{c'-a}$ setiap kali $a<c'\in I$; oleh
karena itu, $\phi(x)\ge\phi(a)+b(x-a)$ untuk setiap $x\in I$.

\medskip

{\bf (b)} Berdasarkan pemilihan $b_q$,
$\phi(x)\ge\sup_{q\in I\cap\Bbb Q}\phi_q(x)$.
Di pihak lain, untuk $x\in I$ yang diberikan, pilih $y\in I$ sedemikian
sehingga $x<y$, lalu ambil $b\in\Bbb R$ sedemikian sehingga
$\phi(z)\ge\phi(x)+b(z-x)$ untuk setiap $z\in I$.
Jika $q\in\Bbb Q$ dan $x<q<y$, kita mempunyai
$\phi(y)\ge\phi(q)+b_q(y-q)$, sehingga
$b_q\le\bover{\phi(y)-\phi(q)}{y-q}$ dan

$$\eqalign{\phi(q)+b_q(x-q)
&=\phi(q)-b_q(q-x)
\ge\phi(q)-\Bover{\phi(y)-\phi(q)}{y-q}(q-x)\cr
&=\Bover{y-x}{y-q}\phi(q)-\Bover{q-x}{y-q}\phi(y)
\ge\Bover{y-x}{y-q}(\phi(x)+b(q-x))-\Bover{q-x}{y-q}\phi(y).\cr}$$

\noindent Sekarang


$$\eqalign{\phi(x)
&=\lim_{q\downarrow x}
   \Bover{y-x}{y-q}(\phi(x)+b(q-x))-\Bover{q-x}{y-q}\phi(y)\cr
&\le\sup_{q\in\Bbb Q\cap\ooint{x,y}}
   \Bover{y-x}{y-q}(\phi(x)+b(q-x))-\Bover{q-x}{y-q}\phi(y)\cr
&\le\sup_{q\in\Bbb Q\cap\ooint{x,y}}\phi(q)+b_q(x-q)
\le\sup_{q\in\Bbb Q\cap I}\phi(q)+b_q(x-q).\cr}$$

\medskip

{\bf (c)} Dengan menuliskan $\phi_q(x)=\phi(q)+b_q(x-q)$ untuk setiap
$q\in\Bbb Q\cap I$, setiap $\phi_q$ merupakan fungsi terukur Borel, dan
$\phi=\sup_{q\in I\cap\Bbb Q}\phi_q$ adalah supremum suatu keluarga
terhitung dari fungsi-fungsi terukur Borel, sehingga terukur Borel.
}%end of proof of 233H

\leader{233I}{Ketaksamaan Jensen} Misalkan $(X,\Sigma,\mu)$ suatu ruang
ukur dan $\phi:\Bbb R\to\Bbb R$ suatu fungsi konveks.

(a) Andaikan $f$ dan $g$ fungsi bernilai real yang terukur secara virtual
terhadap $\mu$ dan didefinisikan hampir di mana-mana dalam $X$, serta
$g\ge 0$ hampir di mana-mana, $\int g=1$, dan $g\times f$
terintegralkan. Maka $\phi(\int g\times f)\le\int g\times\phi f$, dengan
integral di ruas kanan mungkin perlu ditafsirkan sebagai $\infty$.

(b) Khususnya, jika $\mu X=1$ dan $f$ suatu fungsi bernilai real yang
terintegralkan pada $X$, maka $\phi(\int f)\le\int\phi f$.

\proof{{\bf (a)} Untuk setiap $q\in\Bbb Q$, ambil $b_q$ sedemikian
sehingga $\phi(t)\ge\phi_q(t)=\phi(q)+b_q(t-q)$ untuk setiap
$t\in\Bbb R$ (233Ha). Karena $\phi$ terukur Borel (233Hc), $\phi f$
terukur secara virtual terhadap $\mu$ (121H), sehingga $g\times\phi f$
juga demikian; karena $g\times\phi f$ terdefinisi hampir di mana-mana
dan hampir di mana-mana lebih besar daripada atau sama dengan fungsi
terintegralkan $g\times\phi_0f$, $\int g\times\phi f$ terdefinisi dalam
$\ocint{-\infty,\infty}$. Sekarang

$$\eqalign{\phi_q\bigl(\intop g\times f\bigr)
&=\phi(q)+b_q\int g\times f-b_qq\cr
&=\int g\times(b_qf+(\phi(q)-b_qq)\chi X)
=\int g\times\phi_qf
\le\int g\times\phi f,\cr}$$

\noindent karena $\int g=1$ dan $g\ge 0$ hampir di mana-mana. Menurut
233Hb,

\Centerline{$\phi(\intop g\times f)
=\sup_{q\in\Bbb Q}\phi_q(\intop g\times f)
\le\int g\times\phi f$.}

\medskip

{\bf (b)} Ambil $g$ sebagai fungsi konstan bernilai $1$.
}%end of proof of 233I

\leader{233J}{}\cmmnt{ Bahkan kasus khusus 233Ib dari Ketaksamaan Jensen
sudah sangat berguna. Kasus tersebut dapat diperluas sebagai berikut.

\medskip

\noindent}{\bf Teorema} Misalkan $(X,\Sigma,\mu)$ suatu ruang
probabilitas dan $\Tau$ suatu subaljabar-$\sigma$ dari $\Sigma$.
Misalkan $\phi:\Bbb R\to\Bbb R$ suatu fungsi konveks dan $f$ suatu
fungsi bernilai real yang $\mu$-terintegralkan, didefinisikan hampir di
mana-mana dalam $X$, sedemikian sehingga komposisi $\phi f$ juga
terintegralkan. Jika $g$ dan $h$ masing-masing merupakan ekspektasi
bersyarat pada $\Tau$ dari $f$ dan $\phi f$, maka $\phi g\leae h$.
Akibatnya, $\int\phi g\le\int\phi f$.

\proof{ Kita menggunakan gagasan yang sama seperti dalam 233I. Untuk
setiap $q\in\Bbb Q$, ambil $b_q\in\Bbb R$ sedemikian sehingga
$\phi(t)\ge\phi_q(t)=\phi(q)+b_q(t-q)$ untuk setiap $t\in\Bbb R$,
sehingga $\phi(t)=\sup_{q\in\Bbb Q}\phi_q(t)$ untuk setiap
$t\in\Bbb R$. Sekarang tetapkan

\Centerline{$\psi_q(x)=\phi(q)+b_q(g(x)-q)$}

\noindent untuk $x\in\dom g$. Kita melihat bahwa
$\psi_q=\phi_q g$ merupakan ekspektasi bersyarat dari $\phi_q f$, dan
karena $\phi_q f\leae\phi f$, mesti berlaku $\psi_q\leae h$. Tetapi
$\phi g=\sup_{q\in\Bbb Q}\psi_q$ di mana pun $g$ terdefinisi,
sehingga $\phi g\leae h$, sebagaimana dinyatakan.

Langsung diperoleh bahwa $\int\phi g\le\int h=\int\phi f$.
}%end of proof of 233J

\leader{233K}{}\cmmnt{ Saya memberikan proposisi berikut, suatu
pengembangan 233Eg, dalam bentuk yang sangat umum, sebab penerapannya
dapat muncul di mana saja.

\medskip

\noindent}{\bf Proposisi} Misalkan $(X,\Sigma,\mu)$ suatu ruang
probabilitas dan $\Tau$ suatu subaljabar-$\sigma$ dari $\Sigma$.
Andaikan $f$ suatu fungsi yang $\mu$-terintegralkan dan $h$ suatu fungsi
bernilai real yang terukur secara virtual terhadap
$(\mu\restrp\Tau)$ dan didefinisikan
$(\mu\restrp\Tau)$-hampir di mana-mana dalam $X$. Misalkan $g$ dan
$g_0$ ekspektasi bersyarat dari $f$ dan $|f|$ pada $\Tau$. Maka
$f\times h$ terintegralkan jika dan hanya jika $g_0\times h$
terintegralkan, dan dalam hal ini $g\times h$ merupakan ekspektasi
bersyarat dari $f\times h$ pada $\Tau$.

\proof{{\bf (a)} Andaikan $h$ suatu fungsi sederhana terhadap
$\mu\restrp\Tau$. Maka tentu $f\times h$ dan $g_0\times h$
terintegralkan, dan $g\times h$ merupakan ekspektasi bersyarat dari
$f\times h$ seperti dalam 233Eg.

\medskip

{\bf (b)} Sekarang andaikan $f$, $h\ge 0$. Maka $g=g_0\ge 0$ hampir di
mana-mana (233Ec). Misalkan $\tilde h$ suatu fungsi nonnegatif yang
terukur terhadap $\Tau$ dan didefinisikan pada seluruh $X$ sedemikian
sehingga $h\eae\tilde h$. Untuk setiap $n\in\Bbb N$, tetapkan

$$\eqalign{h_n(x)
&=2^{-n}k\text{ jika }0\le k<4^n
  \text{ dan }2^{-n}k\le\tilde h(x)<2^{-n}(k+1),\cr
&=2^n\text{ jika }\tilde h(x)\ge 2^n.\cr}$$

\noindent Maka $h_n$ suatu fungsi sederhana terhadap
$(\mu\restrp\Tau)$, sehingga $g\times h_n$ merupakan ekspektasi
bersyarat dari $f\times h_n$. Barisan $\sequencen{f\times h_n}$ dan
$\sequencen{g\times h_n}$ keduanya merupakan barisan tak menurun hampir
di mana-mana dari fungsi-fungsi terintegralkan, dengan limit masing-masing
$f\times h$ dan $g\times h$. Menurut Teorema B.Levi,

$$\eqalignno{f\times h\text{ terintegralkan}
&\iff\,f\times\tilde h\text{ terintegralkan}\cr
&\iff\,\sup_{n\in\Bbb N}\int f\times h_n<\infty
\iff\,\sup_{n\in\Bbb N}\int g\times h_n<\infty\cr
\noalign{\noindent (karena $\int g\times h_n=\int f\times h_n$ untuk
setiap $n$)}
&\iff\,g\times h\text{ terintegralkan}
\iff\,g_0\times h\text{ terintegralkan}.\cr}$$

\noindent Selain itu, dalam hal ini

$$\eqalign{\int_Ef\times h
&=\int_Ef\times\tilde h
=\lim_{n\to\infty}\int_Ef\times h_n\cr
&=\lim_{n\to\infty}\int_Eg\times h_n
=\int_Eg\times\tilde h
=\int_Eg\times h\cr}$$

\noindent untuk setiap $E\in\Tau$, sedangkan $g\times h$ terukur secara
virtual terhadap $(\mu\restrp\Tau)$, sehingga $g\times h$ merupakan
ekspektasi bersyarat dari $f\times h$.

\medskip

{\bf (c)} Terakhir, tinjau kasus umum dengan $f$ terintegralkan dan $h$
terukur secara virtual. Tetapkan $f^+=f\vee 0$ dan $f^-=(-f)\vee 0$,
sehingga $f=f^+-f^-$ dan $0\le f^+,\,f^-\le|f|$; serupa dengan itu,
tetapkan $h^+=h\vee 0$ dan $h^-=(-h)\vee 0$. Misalkan $g_1$ dan $g_2$
ekspektasi bersyarat dari $f^+$ dan $f^-$ pada $\Tau$. Karena
$0\le f^+$, $f^-\le|f|$, maka $0\le g_1$, $g_2\leae g_0$, sedangkan
$g\eae g_1-g_2$.

Kita melihat bahwa

$$\eqalign{f\times h\text{ terintegralkan}
&\iff\,|f|\times|h|=|f\times h|\text{ terintegralkan}\cr
&\iff\,g_0\times |h|\text{ terintegralkan}\cr
&\iff\,g_0\times h\text{ terintegralkan}.\cr}$$

\noindent Dan dalam hal ini keempat fungsi
$f^+\times h^+,\ldots,f^-\times h^-$ terintegralkan, sehingga

\Centerline{$(g_1-g_2)\times h
=g_1\times h^+-g_2\times h^+-g_1\times h^-+g_2\times h^-$}

\noindent merupakan ekspektasi bersyarat dari

\Centerline{$f^+\times h^+-f^-\times h^+-f^+\times h^-+f^-\times h^- %
=f\times h$.}

\noindent Karena $g\times h\eae(g_1-g_2)\times h$, fungsi ini juga
merupakan ekspektasi bersyarat dari $f\times h$, dan selesailah bukti.
}%end of proof of 233K

\exercises{
\leader{233X}{Latihan dasar (a)}
%\spheader 233Xa
Misalkan $(X,\Sigma,\mu)$ suatu ruang probabilitas dan $\Tau$ suatu
subaljabar-$\sigma$ dari $\Sigma$. Misalkan $\sequencen{f_n}$ suatu
barisan fungsi nonnegatif yang $\mu$-terintegralkan, dan andaikan $g_n$
suatu ekspektasi bersyarat dari $f_n$ pada $\Tau$ untuk setiap $n$.
Andaikan $f=\liminf_{n\to\infty}f_n$ terintegralkan dan mempunyai
ekspektasi bersyarat $g$. Tunjukkan bahwa
$g\leae\liminf_{n\to\infty}g_n$.
%233E

\spheader 233Xb Misalkan $I\subseteq\Bbb R$ suatu selang dan
$\phi:I\to\Bbb R$ suatu fungsi. Tunjukkan bahwa $\phi$ konveks jika dan
hanya jika $\{x:x\in I,\,\phi(x)+bx\le c\}$ merupakan selang untuk setiap
$b$, $c\in\Bbb R$.
%233G

\sqheader 233Xc Misalkan $I\subseteq\Bbb R$ suatu selang terbuka dan
$\phi:I\to\Bbb R$ suatu fungsi. (i) Tunjukkan bahwa jika $\phi$ dapat
didiferensialkan, maka fungsi tersebut konveks jika dan hanya jika
$\phi'$ tak menurun. (ii) Tunjukkan bahwa jika $\phi$ kontinu mutlak
pada setiap subselang tertutup berbatas dari $I$, maka $\phi$ konveks
jika dan hanya jika $\phi'$ tak menurun pada domainnya.

\spheader 233Xd Untuk setiap $r\ge 1$, suatu subhimpunan $C$ dari
$\BbbR^r$ disebut {\bf konveks} jika $tx + (1-t)y\in C$ untuk semua
$x$, $y\in C$, dan $t\in[0,1]$. Jika $C\subseteq\BbbR^r$ konveks, maka
suatu fungsi $\phi:C\to\Bbb R$ disebut {\bf konveks} jika
$\phi(tx+(1-t)y)\le t\phi(x)+(1-t)\phi(y)$ untuk semua $x$, $y\in C$,
dan $t\in[0,1]$.

Misalkan $C\subseteq\BbbR^r$ suatu himpunan konveks dan
$\phi:C\to\Bbb R$ suatu fungsi. Tunjukkan bahwa pernyataan-pernyataan
berikut ekuivalen: (i) fungsi $\phi$ konveks; (ii) himpunan
$\{(x,t):x\in C,\,t\in\Bbb R,\,t\ge\phi(x)\}$ konveks dalam
$\BbbR^{r+1}$; (iii) himpunan
$\{x:x\in C,\,\phi(x)+b\dotproduct x\le c\}$ konveks dalam $\BbbR^r$
untuk setiap $b\in\BbbR^r$ dan $c\in\Bbb R$, dengan menuliskan
$b\dotproduct x=\sum_{i=1}^r\beta_i\xi_i$ jika
$b=(\beta_1,\ldots,\beta_r)$ dan $x=(\xi_1,\ldots,\xi_r)$.
%233G

\spheader 233Xe Misalkan $I\subseteq\Bbb R$ suatu selang dan
$\phi:I\to\Bbb R$ suatu fungsi konveks.

\quad (i) Tunjukkan bahwa jika $a$, $d\in I$, dan $a<b\le c<d$, maka

\Centerline{$\Bover{\phi(b)-\phi(a)}{b-a}
\le\Bover{\phi(d)-\phi(c)}{d-c}$.}

\quad (ii) Tunjukkan bahwa $\phi$ kontinu pada setiap titik interior
$I$.

\quad (iii) Tunjukkan bahwa {\it salah satu} $\phi$ monoton pada $I$
{\it atau} terdapat $c\in I$ sedemikian sehingga
$\phi(c)=\min_{x\in I}\phi(x)$, dan $\phi$ tak menaik pada
$I\cap\ocint{-\infty,c}$ serta monoton tak menurun pada
$I\cap\coint{c,\infty}$.

\quad (iv) Tunjukkan bahwa $\phi$ dapat didiferensialkan di semua titik
$I$ kecuali pada suatu himpunan titik yang paling banyak terhitung, dan turunannya tak
menurun dalam arti bahwa $\phi'(x)\le\phi'(y)$ setiap kali
$x$, $y\in\dom\phi'$ dan $x\le y$.

\quad (v) Tunjukkan bahwa jika $I$ tertutup dan berbatas serta $\phi$
kontinu, maka $\phi$ kontinu mutlak.

\quad (vi) Tunjukkan bahwa jika $I$ tertutup dan berbatas serta
$\psi:I\to\Bbb R$ kontinu mutlak dengan turunan tak menurun, maka $\psi$
konveks.
%233H

\spheader 233Xf Tunjukkan bahwa jika $I\subseteq\Bbb R$ suatu selang dan
$\phi$, $\psi:I\to\Bbb R$ fungsi-fungsi konveks, maka
$a\phi+b\psi$ juga konveks untuk setiap $a$, $b\ge 0$.
%233H

\spheader 233Xg Dalam konteks 233K, berikan contoh ketika $g\times h$
terintegralkan tetapi $f\times h$ tidak terintegralkan.
\Hint{ambil $X$, $\mu$, dan $\Tau$ seperti dalam 233Cb, lalu atur agar
$g$ bernilai $0$.}
%233K

\spheader 233Xh Misalkan $I\subseteq\Bbb R$ suatu selang dan $\Phi$
suatu keluarga tak kosong dari fungsi-fungsi konveks bernilai real pada
$I$ sedemikian sehingga $\psi(x)=\sup_{\phi\in\Phi}\phi(x)$ berhingga
untuk setiap $x\in I$. Tunjukkan bahwa $\psi$ konveks.
%233H

\leader{233Y}{Latihan lanjutan (a)} Jika $I\subseteq\Bbb R$ suatu
selang, suatu fungsi $\phi:I\to\Bbb R$ disebut {\bf konveks-tengah} jika
$\phi(\bover{x+y}2)\le\bover12(\phi(x)+\phi(y))$ untuk semua
$x$, $y\in I$. Tunjukkan bahwa fungsi konveks-tengah yang berbatas pada
suatu subselang tak trivial dari $I$ adalah konveks.
%233G

\spheader 233Yb Perumum 233Xd ke ruang bernorma sembarang sebagai
pengganti $\BbbR^r$.
%233G

\spheader 233Yc Misalkan $(X,\Sigma,\mu)$ suatu ruang probabilitas dan
$\Tau$ suatu subaljabar-$\sigma$ dari $\Sigma$. Misalkan $\phi$ suatu
fungsi konveks bernilai real dengan domain suatu selang
$I\subseteq\Bbb R$, dan $f$ suatu fungsi bernilai real yang terintegralkan
pada $X$ sedemikian sehingga $f(x)\in I$ untuk hampir setiap $x\in X$
dan $\phi f$ terintegralkan. Misalkan $g$ dan $h$ masing-masing
ekspektasi bersyarat pada $\Tau$ dari $f$ dan $\phi f$. Tunjukkan bahwa
$g(x)\in I$ untuk hampir setiap $x$ dan $\phi g\leae h$.
%233J

\spheader 233Yd(i) Tunjukkan bahwa jika $I\subseteq\Bbb R$ suatu selang
berbatas, $E\subseteq I$ terukur Lebesgue, dan
$\mu E>\bover23\mu I$ dengan $\mu$ ukuran Lebesgue, maka untuk setiap
$x\in I$ terdapat $y$, $z\in E$ sedemikian sehingga
$z=\bover{x+y}2$. \Hint{menurut 134Ya/263A,
$\mu(x+E)+\mu(2E)>\mu(2I)$.} (ii) Tunjukkan bahwa jika
$f:[0,1]\to\Bbb R$ suatu fungsi konveks-tengah terukur Lebesgue
(definisi: 233Ya), $a>0$, dan
$E=\{x:x\in[0,1]$, $a\le f(x)<2a\}$ tidak terabaikan, maka terdapat
selang tak trivial $I\subseteq[0,1]$ sedemikian sehingga $f(x)>0$ untuk
setiap $x\in I$. \Hint{223B.} (iii) Andaikan
$f:[0,1]\to\Bbb R$ suatu fungsi konveks-tengah sedemikian sehingga
$f\le 0$ hampir di mana-mana dalam $[0,1]$. Tunjukkan bahwa $f\le 0$ di
setiap titik $\ooint{0,1}$. \Hint{untuk setiap $x\in\ooint{0,1}$,
$\max(f(x-t),f(x+t))\le 0$ untuk hampir setiap
$t\in[0,\min(x,1-x)]$.}
(iv) Andaikan $f:[0,1]\to\Bbb R$ suatu fungsi konveks-tengah terukur
Lebesgue sedemikian sehingga $f(0)=f(1)=0$. Tunjukkan bahwa $f(x)\le 0$
untuk setiap $x\in[0,1]$.
\Hint{tunjukkan bahwa $\{x:f(x)\le 0\}$ padat dalam $[0,1]$, gunakan
(ii) untuk menunjukkan bahwa himpunan tersebut koterabaikan dalam
$[0,1]$, lalu terapkan (iii).}
(v) Tunjukkan bahwa jika $I\subseteq\Bbb R$ suatu selang dan
$f:I\to\Bbb R$ suatu fungsi konveks-tengah terukur Lebesgue, maka fungsi
tersebut konveks.
%233Ya 233G not in 2001 edition

\spheader 233Ye Misalkan $(X,\Sigma,\mu)$ suatu ruang probabilitas,
$\Tau$ suatu subaljabar-$\sigma$ dari subhimpunan-subhimpunan $X$, dan
$f:X\to[0,\infty]$ suatu fungsi terukur terhadap $\Sigma$. Tunjukkan
bahwa (i) terdapat suatu fungsi yang terukur terhadap $\Tau$, yaitu
$g:X\to[0,\infty]$, sedemikian sehingga $\int_Fg=\int_Ff$ untuk setiap
$F\in\Tau$
(ii) setiap dua fungsi semacam itu sama hampir di mana-mana.

\spheader 233Yf Andaikan $r\ge 1$ dan
$C\subseteq\BbbR^r\setminus\{0\}$ suatu himpunan konveks. Tunjukkan
bahwa terdapat $b\in\BbbR^r$ tak nol sedemikian sehingga
$\varinnerprod{b}{z}\ge 0$ untuk setiap $z\in C$.
\Hint{jika $r=2$, identifikasikan $\BbbR^2$ dengan $\Bbb C$; reduksikan
ke kasus ketika $C$ tidak memuat titik yang real dan negatif; tetapkan
$\theta=\sup\{\arg z:z\in C\}$ dan $b=-ie^{i\theta}$. Lalu gunakan
induksi pada $r$.}
%233Xd

\spheader 233Yg Andaikan $r\ge 1$, $C\subseteq\BbbR^r$ suatu himpunan
konveks, dan $\phi:C\to\Bbb R$ suatu fungsi konveks. Tunjukkan bahwa
terdapat fungsi $h:\BbbR^r\to\coint{-\infty,\infty}$ sedemikian sehingga
$\phi(z)=\sup\{h(y)+\varinnerprod{z}{y}:y\in\BbbR^r\}$ untuk setiap
$z\in C$.
\Hint{coba $h(y)=\inf\{\phi(z)-\varinnerprod{z}{y}:z\in C\}$, lalu
terapkan 233Yf pada suatu translasi dari
$\{(z,t):\phi(z)\le t\}$.}
%233Yf 233Xd 233H

\spheader 233Yh Misalkan $(X,\Sigma,\mu)$ suatu ruang probabilitas,
$r\ge 1$ suatu bilangan bulat, dan $C\subseteq\BbbR^r$ suatu himpunan
konveks. Misalkan $f_1,\ldots,f_r$ fungsi-fungsi bernilai real yang
$\mu$-terintegralkan dan andaikan
$\{x:x\in\bigcap_{j\le r}\dom f_j$,
$(f_1(x),\ldots,f_r(x))\in C\}$ merupakan subhimpunan koterabaikan dari
$X$. Tunjukkan bahwa $(\int f_1,\ldots,\int f_r)\in C$.
\Hint{gunakan induksi pada $r$.}
%233Yg 233H

\spheader 233Yi Misalkan $(X,\Sigma,\mu)$ suatu ruang probabilitas,
$r\ge 1$ suatu bilangan bulat, $C\subseteq\BbbR^r$ suatu himpunan
konveks, dan $\phi:C\to\Bbb R$ suatu fungsi konveks. Misalkan
$f_1,\ldots,f_r$ fungsi-fungsi bernilai real yang
$\mu$-terintegralkan dan andaikan
$\{x:x\in\bigcap_{j\le r}\dom f_j$,
$(f_1(x),\ldots,f_r(x))\in C\}$ merupakan subhimpunan koterabaikan dari
$X$. Tunjukkan bahwa
$\phi(\int f_1,\ldots,\int f_r)\le\int\phi(f_1,\ldots,f_r)$.
%233Yg 233Yh 233H

\spheader 233Yj Misalkan $(X,\Sigma,\mu)$ suatu ruang ukur dengan
$\mu X>0$, $r\ge 1$ suatu bilangan bulat, $C\subseteq\BbbR^r$ suatu
himpunan konveks sedemikian sehingga $tz\in C$ setiap kali $z\in C$ dan
$t>0$, dan $\phi:C\to\Bbb R$ suatu fungsi konveks. Misalkan
$f_1,\ldots,f_r$ fungsi-fungsi bernilai real yang
$\mu$-terintegralkan dan andaikan
$\{x:x\in\bigcap_{j\le r}\dom f_j$,
$(f_1(x),\ldots,f_r(x))\in C\}$ merupakan subhimpunan koterabaikan dari
$X$. Tunjukkan bahwa $(\int f_1,\ldots,\int f_r)\in C$ dan
$\phi(\int f_1,\ldots,\int f_r)\le\int\phi(f_1,\ldots,f_r)$.
\Hint{dengan menggabungkan 215B(viii) dan 235K di bawah, tunjukkan bahwa
terdapat suatu ukuran probabilitas $\nu$ pada $X$ dan suatu fungsi
$h:X\to\coint{0,\infty}$ sedemikian sehingga
$\int f_jd\mu=\int f_j\times h\,d\nu$ untuk setiap $j$.}
%233Yh 233Yi 233H

}%end of exercises

\endnotes{
\Notesheader{233} Konsep `ekspektasi bersyarat` bersifat fundamental
dalam teori probabilitas dan akan muncul kembali dalam konteks
alamiahnya pada Bab 27. Namun, bahkan sebagai latihan teknik, saya
berharap konsep tersebut akan terasa patut diperhatikan.

Saya memperkenalkan 233E sebagai `daftar fakta elementer`, dan
fakta-fakta tersebut memang langsung. Namun, di bawah permukaan terdapat
sejumlah gagasan luar biasa yang menunggu untuk diungkapkan. Jika Anda
mengambil $\Tau$ sebagai aljabar trivial $\{\emptyset,X\}$, sehingga
ekspektasi bersyarat (yang unik) dari suatu fungsi terintegralkan $f$
adalah fungsi konstan $(\int f)\chi X$, maka 233Ed dan 233Ee menjadi
versi-versi Teorema B.Levi dan Teorema Konvergensi Terdominasi Lebesgue.
(Lema Fatou terdapat dalam 233Xa.) Bahkan 233Eg dapat dipandang sebagai
perumuman hasil $\int cf=c\int f$, dengan konstanta $c$ diganti oleh
fungsi berbatas yang terukur terhadap $\Tau$. Tema berulang dalam
bagian-bagian selanjutnya risalah ini ialah pencarian versi `bersyarat`
dari teorema-teorema. Bukti 233Ee merupakan contoh khas argumen yang
diterjemahkan dari bukti hasil `tak bersyarat` semula.

Saya mengemukakan bahwa 233I-233J mengejutkan, dan menurut saya sebagian
besar dari kita memang merasakannya demikian, bahkan ketika hasil-hasil
itu diterapkan pada daftar fungsi konveks yang diberikan dalam 233G.
Namun, perlu saya nyatakan bahwa dalam satu segi 233J hanya sedikit
berkaitan dengan ekspektasi bersyarat. Satu-satunya sifat ekspektasi
bersyarat yang digunakan dalam bukti ialah (i) jika $g$ suatu ekspektasi
bersyarat dari $f$, maka $a\chi X+bg$ merupakan ekspektasi bersyarat
dari $a\chi X+bf$ untuk semua $a$, $b$ real (ii) jika $g_1$, $g_2$
merupakan ekspektasi bersyarat dari $f_1$, $f_2$ dan $f_1\leae f_2$,
maka $g_1\leae g_2$. Lihat 244Xm di bawah. Ketaksamaan Jensen mempunyai
perluasan menarik ke kasus multidimensi, yang ditelaah dalam
233Yf-233Yj.   %233Yf 233Yg 233Yh 233Yi 233Yj
Jika Anda pernah menjumpai bentuk `geometris` Teorema Hahn-Banach (lihat
3A5C dalam Jilid 3), Anda akan mendapati 233Yf dan 233Yg sangat alami,
dan mungkin memperhatikan bahwa kasus berdimensi hingga agak berbeda
dari kasus berdimensi tak hingga yang mungkin pernah diajarkan kepada
Anda. Menurut saya, langkah yang sebenarnya paling pelik terdapat dalam
233Yh.

Perhatikan bahwa 233Ib dapat dipandang sebagai kasus khusus 233J ketika
$\Tau=\{\emptyset,X\}$. Bahkan 233Ia dapat diturunkan dari 233Ib yang
diterapkan pada ukuran $\nu$ dengan $\nu E=\int_Eg$ untuk setiap
$E\in\Sigma$.

Seperti 233B, argumen dalam 233K tampaknya memiliki cukup banyak
perincian teknis. Pokok hasil ini ialah bahwa kita dapat menyimpulkan
keterintegralan $f\times h$ dari keterintegralan $g_0\times h$ (tetapi
tidak dari keterintegralan $g\times h$; lihat 233Xg).
Selebihnya seharusnya rutin.
}%end of notes

\discrpage
